\section*{Results}

\subsection*{Relational vs flat annotation priors}

\textbf{Flat-prior fine-mapping collapses relational structure to a
scalar.} In the SuSiE/FINEMAP/PolyFun family, biological annotations
enter as a per-variant weight vector $w \in \mathbb{R}^n$, averaged
across tissues and detached from the gene, pathway and interaction
context each variant belongs to. GraphGWAS instead represents this
context as a typed graph (Figure~\ref{fig:graph}): variants connect
to genes through $\texttt{HAS\_CONSEQUENCE}$ and tissue-specific
$\texttt{eQTL}$ edges, genes participate in pathways via
$\texttt{IN\_PATHWAY}$, and the gene layer is coupled within itself
by protein--protein $\texttt{INTERACTS\_WITH}$ edges. The graph loaded on the 1000 Genomes Phase~3 high-coverage cohort
\cite{1000genomes2015global,byrska2022high} carries 70.7~million
variants, 20{,}092 GENCODE~v47 genes\cite{frankish2023gencode},
43.2~million GTEx~v8 tissue eQTLs\cite{gtex2020}, 230{,}850 STRING
interactions at combined score~$\geq$~700\cite{szklarczyk2023string},
and 370{,}000 ENCODE cCRE regulatory elements\cite{encode2020expanded}
(full counts in Supplementary Table~S1). All
relational edges are first-class: a fine-mapping procedure can either
collapse them to a scalar prior (the flat-prior approach) or carry
them through the inference (the relational approach we develop
below).

\begin{figure}[htbp]
\centering
\includegraphics[width=0.95\textwidth]{fig1_architecture.pdf}
\caption{\textbf{Relational vs flat annotation priors in fine-mapping.}
(\textbf{a})~Flat-prior fine-mappers (PolyFun / SuSiE family) accept
a vector of per-variant annotation weights $w_i$ that collapses
gene, tissue, pathway and interaction information into a single
scalar per variant; the Bayesian regression proceeds on
($z$-scores, LD matrix $R$, prior weights $w$) and returns a PIP
vector whose biological context must be reattached by a post-hoc
enrichment step.  (\textbf{b})~GraphGWAS exposes the multi-omics
knowledge graph directly: a query variant (highlighted red circle)
is connected by \emph{typed} edges---\texttt{HAS\_CONSEQUENCE} to
its coding genes, tissue-specific \texttt{eQTL} edges to the genes
it regulates, \texttt{IN\_REGULATORY} to ENCODE cCREs,
\texttt{IN\_PATHWAY} via genes to pathway nodes, and STRING
\texttt{INTERACTS\_WITH} edges coupling genes at combined
score~$\geq$~700---to five typed node classes (Variant, Gene,
Pathway, RegulatoryElement, Tissue; legend bottom-left). Real
labels are illustrative (FTO, IRX3, Adipose\_Subcutaneous, fatty-
acid metabolism). The credible-set output is itself a graph object:
each reported variant is co-queryable with its gene, tissue and
pathway neighbours in one traversal. The HBP / L1 algorithm that
computes posterior inclusion probabilities by message passing
across this graph is detailed in Figure~\ref{fig:hbp}; this
panel shows the underlying data model, not the inference
procedure. Multi-omics graph state on 1000 Genomes Phase~3 (inline
lower-right box; full counts in Supplementary Table~S1): 70.7~M
variants, 20{,}092 GENCODE~v47 genes, 43.2~M GTEx~v8 eQTL edges
across 49 tissues, 230{,}850 STRING interactions, and 370{,}000
ENCODE cCRE regulatory elements.}
\label{fig:graph}
\end{figure}

\subsection*{Hierarchical belief propagation on a three-layer factor graph}

HBP converts the graph into iterated message passing across
variant, gene and pathway layers (Figure~\ref{fig:hbp}). Each
iteration runs an upward pass---variant beliefs aggregate to gene
scores weighted by eQTL strength, then diffuse across the gene layer
along STRING PPI edges at a 0.3 coupling, then aggregate to
pathway scores---followed by a downward pass that returns a variant
prior $\pi^{(t)}$ via the same edges. The variant belief is combined
as $b^{(t+1)} = \alpha \cdot \text{softmax}(u) + (1-\alpha) \pi^{(t)}$,
where $u$ is the LD-deconvolved z-statistic, and damped by
$\lambda$ before renormalisation (defaults $\alpha = 0.6$,
$\lambda = 0.5$, 5 rounds). Under positive damping the iteration is
a strict contraction on the probability simplex, so by Banach's
fixed-point theorem HBP converges geometrically to a unique
self-consistent belief (Theorem~2, Supplementary Note~S1). At
default hyperparameters the algorithm reaches four-decimal
convergence in $\approx$0.08~s per locus on 1000 Genomes chromosome~22,
approximately 20--30$\times$ faster than SuSiE and FINEMAP on the
same data.

\begin{figure}[htbp]
\centering
\includegraphics[width=0.85\textwidth]{fig2_hbp_schematic.pdf}
\caption{\textbf{Hierarchical belief propagation on a three-layer factor graph.}
HBP performs message passing on a typed factor graph with variant
nodes (bottom), gene nodes (middle), and pathway nodes (top). The
upward pass aggregates variant beliefs into gene scores weighted by
eQTL strength, diffuses them once along the STRING
protein--protein interaction graph at a fixed 0.3 coupling, and
aggregates to the pathway layer. The downward pass reverses the
direction and produces a variant prior $\pi^{(t)}$. Each iteration
combines the prior with LD-deconvolved statistical evidence as
$b^{(t+1)} = \alpha \cdot \text{softmax}(u) + (1-\alpha) \pi^{(t)}$
and damps by $\lambda$ before renormalisation. Defaults are
$\alpha = 0.6$, $\lambda = 0.5$, $T = 5$ rounds. Under positive
damping the update is a strict $\ell_1$ contraction on the
probability simplex (Theorem~2, Supplementary Note~S1), guaranteeing
geometric convergence to a unique fixed point. See
Methods~\ref{methods:hbp} for the full algorithm.}
\label{fig:hbp}
\end{figure}

\subsection*{L1 Bayesian fine-mapping wins weak signal 27--2 against SuSiE}

The regime that matters most for post-GWAS discovery is the one
where statistical signal alone is insufficient to resolve the causal
variant: sub-genome-wide-significant loci, rare-variant effects, and
complex LD blocks. We designed L1 Bayesian fine-mapping for this
regime. L1 first LD-deconvolves the z-statistics
($u_i = z_i - \frac{1}{|N(i)|} \sum_{j \in N(i)} r^2_{ij} z_j$ with
$N(i) = \{j : r^2_{ij} > 0.3\}$), then combines the deconvolved
evidence with a multi-layer graph functional score $z_{\text{func}}$
via $s_i = \alpha u_i + (1-\alpha) z_{\text{func},i}$ with
$\alpha$ chosen by empirical Bayes on a held-out calibration set.
Under mild LD-decay assumptions the causal variant has the highest
unique-signal contribution (Theorem~3).

On 79 independent simulated loci with tissue-specific eQTL causal
variants and weak effect size ($\beta = 0.15$, $h^2 = 0.01$), L1
resolves 61/79 loci at rank~\#1 versus SuSiE's 45/79
(Figure~\ref{fig:weak_signal}). Mean causal-variant rank is 1.65
for L1 versus 2.84 for SuSiE. Head-to-head on matched replicates, L1
beats SuSiE on 27, ties on 50, and loses on 2 (win ratio 13.5\,:\,1,
sign-test $p < 10^{-5}$). The runtime advantage persists: 0.07~s per
locus for L1 versus 1.8~s for SuSiE.

The 27--2 headline is regime-specific, not universal. A 100-replicate
replication on 1000 Genomes chromosome~22 without tissue-specific
eQTL causal enrichment---$\beta = 0.2$, $h^2 = 0.02$, 30 rotating
centres across 17--46~Mb---converges to parity between L1 (54\%
rank-\#1) and SuSiE (57\%), with L1 retaining a $\sim$25$\times$
runtime advantage. The two results together bound the claim:
\textbf{L1's decisive statistical advantage emerges when annotation
informativeness is high} (tissue-specific eQTLs, moderate-impact
functional consequences, pathway co-enrichment); \textbf{under
uniform annotation conditions the advantage narrows to parity on
rank and persists only on speed.} This is the pattern that a
relational prior should produce, not a universally stronger
statistical test.

\begin{figure}[htbp]
\centering
\includegraphics[width=0.95\textwidth]{fig6_weak_signal_headline.pdf}
\caption{\textbf{L1 beats SuSiE 27--2 head-to-head at weak signal with
tissue-specific eQTL priors.} 79 independent simulated loci on
1000 Genomes chromosome~22, with causal variants drawn from GTEx~v8
significant eQTLs at tissue-specificity $-\log_{10} p > 15$ in at
most two tissues; effect size $\beta = 0.15$, $h^2 = 0.01$.
(\textbf{a})~Distribution of causal-variant rank (lower is better):
L1 mean 1.65, SuSiE mean 2.84. (\textbf{b})~Head-to-head scoreboard:
L1 wins 27, ties 50, SuSiE wins 2; win ratio 13.5\,:\,1
(sign-test $p < 10^{-5}$). (\textbf{c})~Rank-\#1 rate: L1~77\%
(61/79) vs SuSiE~57\% (45/79). (\textbf{d})~Per-replicate scatter
of L1 rank against SuSiE rank. Points above the diagonal are
SuSiE-favoured replicates (2/79); points below are L1-favoured
(27/79). Runtime 0.07~s for L1 vs 1.8~s for SuSiE. At higher
annotation uniformity (see next subsection) the L1 advantage
narrows to parity on rank.}
\label{fig:weak_signal}
\end{figure}

\subsection*{Head-to-head against state-of-the-art baselines}

A common GraphGWAS interface integrates the six methods benchmarked
by Wu~\emph{et~al.}~2026\cite{wu2026genome}: SuSiE\cite{wang2020simple},
FINEMAP\cite{benner2016finemap}, SuSiE-inf\cite{cui2024improving},
FINEMAP-inf\cite{cui2024improving}, a Polyfun-proxy\cite{weissbrod2020polyfun}
that feeds the same eQTL priors to SuSiE via
\texttt{prior\_weights}, and SBayesRC\cite{zheng2024leveraging}. On 30
1000 Genomes chromosome~22 F1 simulations at $h^2 = 0.05$ (Figure~\ref{fig:baselines};
Supplementary Table~S2), SuSiE-inf marginally improves on SuSiE
(22/30 vs 21/30 rank-\#1), consistent with Wu~\emph{et~al.}; HBP
reaches 20/30 rank-\#1 at 6--60$\times$ the speed of every baseline
(HBP 0.016~s; SuSiE-inf 0.16~s; FINEMAP-inf 0.20~s; SuSiE 0.98~s;
FINEMAP 2.2~s). When the eQTL prior is also supplied to SuSiE via
Polyfun-proxy on a parallel weak-signal benchmark, SuSiE's rank-\#1
count rises from 10/20 to 11/20 and L1 matches at 26$\times$ the
speed. Combined with the 27--2 weak-signal result, the head-to-head
landscape resolves: \textbf{graph-native and flat-vector priors
converge on accuracy when the prior information is identical; the
graph wins decisively on speed, and on accuracy in regimes where
tissue-specific eQTL structure is informative.}

The 6--60$\times$ speed advantage is partly statistical (fewer
latent effects to fit) and partly a property of the data model.
Reconstructing each locus's multi-omics neighbourhood from
separate files (BED tracks for regulatory elements, GTF for genes,
TSV matrices for tissue-specific eQTLs) is index-bound;
the same information stored as a typed graph collapses the join
into a single index-free traversal that returns in milliseconds.
SBayesRC occupies a complementary problem (genome-wide multi-component
Bayesian mixture on $\sim$1.2~M HapMap~3 SNPs requiring $N \gtrsim
10^5$); HBP and L1 target single-locus resolution at leads surfaced
by genome-wide methods.

\begin{figure}[htbp]
\centering
\includegraphics[width=0.95\textwidth]{fig3_hbp_vs_susie_finemap.pdf}
\caption{\textbf{HBP and L1 vs state-of-the-art Bayesian baselines.}
Head-to-head benchmark on 90 simulated loci from 1000 Genomes
chromosome~22, spanning three scenarios (strong signal; weak signal;
functional-causal with eQTL priors). HBP matches SuSiE rank-\#1
rate within 3 percentage points across all scenarios (70\% strong,
60\% weak, 37\% functional). Runtime distributions: HBP~0.016--0.08~s,
L1~0.055--0.07~s, SuSiE-inf~0.16~s, FINEMAP-inf~0.20~s, SuSiE~1.0--1.8~s,
FINEMAP~2.2~s. When the eQTL prior used by HBP/L1 is also supplied
to SuSiE via the Polyfun-proxy mechanism, SuSiE's rank-\#1 rate
improves from 10/20 to 11/20 but remains 26$\times$ slower than L1
at parity accuracy. (\textbf{a})~Rank distributions; (\textbf{b})~
head-to-head scoreboard; (\textbf{c})~rank-\#1 rate across scenarios;
(\textbf{d})~per-replicate scatter. Full numerical breakdown in
Supplementary Table~S2.}
\label{fig:baselines}
\end{figure}

\subsection*{Calibration, null false-positive rate, and sample-size scaling}

PIPs must mean what they say and must not fire under the null. On
100 null simulations (no causal variant, Gaussian phenotype) over
random 50-kb yeast loci, every method (HBP, L1, SuSiE, FINEMAP)
reports $\max$\,PIP~$<$~0.5 in every replicate, with mean
$\max$\,PIP at 0.003--0.029 --- matching to within 2\% the
softmax bound $\mathbb{E}[\max_i \pi_i] \leq
O(\sqrt{\log n}/n)$ proved in Supplementary Note~S1. On 200
calibration simulations across $h^2 \in \{0.02, 0.05, 0.10, 0.20\}$,
true-discovery rates track the reported PIP diagonal for SuSiE,
FINEMAP and L1; HBP is mildly conservative at mid-range PIPs but
when it assigns PIP $>$ 0.5 the variant is causal 99\% of the
time (Supplementary Figure~S4).

HBP fine-mapping quality scales monotonically with sample size:
on 1000 Genomes chromosome~22 subsampled at
$N \in \{500, 1{,}000, 2{,}000, 3{,}000\}$ ($h^2 = 0.10$, causal
MAF 10--40\%, 30 replicates per $N$), mean causal PIP rises
0.54~$\rightarrow$~0.58~$\rightarrow$~0.59~$\rightarrow$~0.77, with
73\% of replicates at rank~1 by $N = 3{,}000$ (Supplementary
Figure~S5). The linear trend extrapolates beyond the training
regime, supporting the biobank-scale claim tested directly on
Pan-UK Biobank in the next subsection.

\subsection*{Cross-ancestry fine-mapping: 1KG simulated and Pan-UKB real data}\label{sec:cross_anc}

LD structure differs substantially across ancestries and many causal
variants are ancestry-private. On 1KG chromosome~22 restricted to
the three largest superpopulations (EUR $n=503$, AFR $n=661$, EAS
$n=504$, F1 simulations at $h^2 = 0.05$, 30 replicates each), HBP
rank-\#1 ordering is \textbf{AFR 53\% $>$ EAS 27\% $>$ EUR 17\%}
(mean PIPs 0.35, 0.21, 0.08; Figure~\ref{fig:cross_anc}a) --- the
expected direction, since AFR's older, less-correlated LD produces
fewer indistinguishable proxies per locus. HBP's precision holds
across all three ancestries.

A sumstats-only entry path consumes Pan-UK
Biobank\cite{karczewski2024pan,bycroft2018uk} summary statistics
streamed via \texttt{tabix}\cite{li2011tabix} over HTTPS, with
GRCh37$\to$GRCh38 liftover to align with the 1KG reference panel.
Applied to four canonical multi-ancestry loci ---
\emph{FTO}/BMI\cite{frayling2007common},
\emph{APOA5}/triglycerides\cite{do2013common},
\emph{LDLR}/LDL-C\cite{teslovich2010biological},
\emph{HMGA2}/height\cite{weedon2008genome,yengo2022saturated} ---
across four Pan-UKB ancestries (EUR $N=420{,}531$; CSA $8{,}876$;
AFR $6{,}636$; EAS $2{,}709$), three of the four EUR cells (BMI,
LDL, height) resolve to a \textbf{single-variant credible set at
PIP $=$ 1.000}. \emph{HMGA2}/height is cross-ancestry robust: CSA
PIP $=$ 1.000 at $N=8{,}876$ and EAS PIP $=$ 0.97 at $N=2{,}709$ ---
a 145$\times$-smaller sample than EUR. \emph{APOA5}/triglycerides
shows two-ancestry convergence on the same indel
(EUR PIP $=$ 0.51, CSA PIP $=$ 0.73). The same codebase and the
same algorithm; no EUR-specific tuning is needed.

\begin{figure}[htbp]
\centering
\includegraphics[width=0.95\textwidth]{fig9_cross_ancestry.pdf}
\caption{\textbf{Cross-ancestry fine-mapping: controlled 1KG
simulation and real Pan-UK Biobank demonstration.}
(\textbf{a})~Controlled simulation: 1KG chromosome~22 restricted to
three superpopulations (EUR $n{=}503$, AFR $n{=}661$, EAS $n{=}504$);
F1 simulations at $h^2 = 0.05$, 30 replicates per ancestry. HBP
rank-\#1 rate AFR~53\% $>$ EAS~27\% $>$ EUR~17\% reflects AFR's older,
less-correlated LD. (\textbf{b})~Real-biobank demonstration: HBP
sumstats-only fine-mapping on four canonical loci (FTO~/~BMI,
APOA5~/~triglycerides, LDLR~/~LDL-C, HMGA2~/~height) across four
Pan-UKB ancestries (EUR $N{=}420{,}531$; CSA $8{,}876$; AFR $6{,}636$;
EAS $2{,}709$). Three EUR cells resolve to a single-variant 95\%
credible set at PIP~=~1.0. Height/HMGA2 is cross-ancestry robust:
CSA PIP~=~1.0, EAS PIP~=~0.97 at 145$\times$ smaller $N$ than EUR.
Triglycerides/APOA5 shows two-ancestry convergence on the same indel.
No EUR-specific tuning; no algorithmic change between the simulated
and real-biobank panels.}
\label{fig:cross_anc}
\end{figure}

\subsection*{Cross-species fine-mapping at biobank and germplasm-bank scale}
\label{sec:multispecies}

Fine-mapping causal variants is hard for distinct reasons across the
scales human medical genetics, model-organism crosses, and crop
germplasm operate at: in human cohorts, signal-to-noise is high but
LD is dense and reference panels are ancestry-specific; in
inbred-yeast and \emph{Arabidopsis} panels, LD is exceptionally long
and pathway/PPI annotation is sparse; in crop germplasm such as the
3{,}000 Rice Genomes panel, accessions span large between-population
$F_{ST}$ and most traits have only ordinal phenotypic scoring with
no curated causal-gene catalogue. The same fine-mapping pipeline,
unchanged except for the loaded annotation graph, was applied to all
three regimes (Supplementary Tables~S5--S6).

\paragraph{Rice (3{,}000 Rice Genomes, 18 morphological traits).}
A whole-genome GWAS on 3{,}024 accessions (29.6~M biallelic SNPs)
across 18 IRRI Standard Evaluation System traits\cite{wang2018genomic}
($\lambda_{GC}$: 0.50--2.07, median 1.09; per-trait detail in
Supplementary Table~S5) provided the GWAS substrate for fine-mapping
the top three genome-wide-significant plus one suggestive
independent lead per trait (72 total) within $\pm$250~kb windows
against a rice multi-omics graph (RAP-DB variant--gene + Ren 2023
gene--pathway curation\cite{ren2023rice} + RicePPINet
gene--gene\cite{liu2017riceppinet}).

L1 delivers a single-variant 95\% credible set with PIP~$\geq$~0.99
on 11 of 72 leads (15\%). \textbf{Twenty-four of 72 leads (33\%)
fine-map to within 250~kb of a Ren-2023 catalogued causal gene}, with
14 of 18 traits contributing at least one validated locus. Notable
hits include the rice red-pericarp gene \emph{Rc} only \textbf{4.5~kb}
from the spikelet-pericarp-colour lead at Chr7:6.07~Mb;
\textbf{\emph{TGW6}} 12.3~kb from the leaf-length-type Chr6:25.11~Mb
lead; the awn-development gene \emph{GAD1/RAE2} 4.4~kb from the
awning-presence Chr8:23.98~Mb lead ($p=4.4\times10^{-90}$,
PIP~$=$~0.999); grain-shape genes \textbf{\emph{An-1}}
($240$~kb, blade-colour Chr4:16.47~Mb, PIP~$=$~1.00) and
\textbf{\emph{BG1}} ($68.7$~kb, spikelet-fertility Chr3:4.10~Mb);
the panicle-development gene \emph{OsER1} ($62$~kb,
apiculus-colour Chr6:5.31~Mb, PIP~$=$~0.77); and the seed-dormancy
regulator \emph{OsVP1} ($176$~kb, panicle-type Chr1:39.54~Mb).
The two highest-confidence pigment-related leads validate against
the same axis of variation captured by the IRRI ordinal scoring
(spikelet pericarp at Chr7:6.07~Mb / \emph{Rc} and apiculus colour
at Chr10:13.32~Mb / \emph{BRD2/LTBSG1}, $4.5$~kb and $28.5$~kb
respectively, both at PIP $\geq 0.99$). Four traits with
$\lambda_{GC} \in [1.6, 2.1]$ show residual subpopulation
confounding ($F_{ST} \approx 0.5$ between XI and GJ subpopulations
under 10-PC correction); the eight doubly-resolved loci with
HBP CS $=$ L1 CS $=$ 1 are all in well-calibrated traits
($\lambda_{GC} \in [0.50, 1.38]$).

\paragraph{Yeast (1011 Genome panel, 35 growth traits).}
Across 35 grammar-corrected growth-trait sumstats\cite{peter2018genome,
bloom2015genetic} fine-mapped against an SGD-derived
gene--GO-slim graph (1.14~M variants, 60\% of the panel), 245 leads
were fine-mapped at $\pm$15~kb windows. Yeast LD is strong enough that
no locus reaches L1 CS $=$ 1, but \textbf{HBP narrows the credible
set on 244 of 245 leads (100\%)}, with 5 loci reaching PIP $\geq 0.5$.
Forty-four of the 245 leads (18\%) lie within 250~kb of a
literature-curated trait-relevant gene (e.g.\ TUP1/SPT7 for ethanol,
CUP1/CUP2 for copper resistance, the GAL pathway for galactose,
ENA1/HOG1 for salt response, ERG/PDR for fluconazole).

\paragraph{Arabidopsis (1001 Genomes, 14 priority traits).}
A TAIR10 + Plant Reactome + STRING-\emph{Ath} graph (covering
1.06--3.99~M variants across the five chromosomes; 20{,}192 genes
with high-confidence PPI partners after UniProt-orthologue
resolution) supports a multi-trait \texttt{plink2} GWAS on 14
priority phenotypes (11 AraGWAS-Bonferroni-validated traits +
flowering-time traits FT10, FT16 and FLC). Of the 54 fine-mapped
leads, three reach L1 CS $=$ 1 with PIP $\geq 0.98$ (a seed-storage
trait at three independent chr4 and chr5 loci); HBP narrows the
credible set on 29 of 54 (54\%). Validation against AraGWAS plus
canonical flowering-time genes (FRI, FLC, CONSTANS, GIGANTEA, VIN3,
FT, DOG1, HKT1) gives 3 of 54 (6\%); the apparently low rate
reflects the chr-4-only coverage of the AraGWAS Bonferroni file
rather than poor fine-mapping.

\paragraph{Human (Pan-UKB EUR, four phenotypes, all 22 autosomes).}
A genome-wide GTEx + STRING graph (\textbf{4.59~M annotated
variants}; 49 GTEx tissues + STRING v12 PPI at score $\geq$ 700;
GENCODE v47 gene model) supports tabix-over-HTTPS Pan-UKB summary
statistics with on-the-fly GRCh37$\to$GRCh38 liftover and per-chromosome
1KG NYGC LD references (3{,}202 samples, GRCh38). \textbf{Three
hundred and twenty-one leads were fine-mapped across BMI, LDL
direct, and triglycerides}; the Pan-UKB Height sumstats are heavily
flagged \texttt{low\_confidence} genome-wide under the
500~kb-clustering filter we applied and returned no leads per
chromosome (the locus-targeted Height/HMGA2 query reported in
\S\ref{sec:cross_anc} bypasses this clustering by querying a specific
500~kb window directly).

\textbf{Genome-wide L1 fine-mapping recovered the textbook canonical
lipid-metabolism genes at single-variant resolution, without any
annotation prior:} \textbf{\emph{LDLR}} (LDL chr19:11.08~Mb, L1
CS $=$ 1, PIP $=$ 1.00, 12.7~kb from the Brown \& Goldstein 1986
prototype gene); \textbf{\emph{APOE}} (LDL and triglycerides
chr19:44.81/44.91~Mb, PIP $=$ 0.60--0.77);
\textbf{\emph{LPL}} (chr8:20.0~Mb, L1 CS $=$ 1, PIP $=$ 0.9998);
\textbf{\emph{GCKR}} (chr2:27.5~Mb, L1 CS $=$ 2, PIP $=$ 0.92,
\textbf{0~bp} from the gene body); \textbf{\emph{ANGPTL3}}
(chr1:62.97~Mb, PIP $=$ 0.81, \textbf{0~bp});
\textbf{\emph{ANGPTL4}} (chr19:8.4~Mb); and the
\textbf{\emph{APOA5} cluster} (chr11:117.0~Mb, PIP $=$ 0.65). Eight of
321 leads reach L1 CS $=$ 1; 47 reach PIP $\geq 0.5$;
the strong-PIP subset (L1 CS $\leq 5$ \emph{and} PIP $\geq 0.5$)
matches a curated chr-22 / GWAS-Catalog literature catalogue at 9 of
10 leads. Across the genome-wide scan, HBP and L1 produce identical
credible sets on every one of 321 leads --- a finding to which we
return below.

\subsection*{Heterogeneity, not density, makes graph priors work}
\label{sec:heterogeneity}

The cross-species data above show 287 of 692 leads (41\%)
fine-mapped tighter under HBP than under flat-prior L1, but the
distribution of that gain across species is striking and
non-uniform: yeast 100\%, Arabidopsis 54\%, rice 19\%, human 0\%
genome-wide. The naive interpretation that "more annotation density
helps fine-mapping" is therefore \emph{wrong} as a general rule:
the human cache is by far the densest of the four (4.59~M
annotated variants, 49 tissues, near-saturating GTEx + STRING
coverage) yet shows zero HBP narrowing relative to L1. The
explanation is that the human cache is also nearly \emph{uniform}
per variant --- almost every annotated variant has the same
neighbourhood structure (one eGene plus tens of STRING partners).
Random graph priors with low per-variant variation cannot break
LD ties that flat-prior fine-mapping has not already broken.

\paragraph{Ablation experiment.} We tested this directly on the
3{,}000 Rice Genomes panel by randomly subsampling the rice
multi-omics cache at five retention fractions (10\%, 25\%, 50\%, 75\%,
100\%) $\times$ 10 bootstrap seeds, and re-running HBP on each of
five simulated grain-quality causal loci with curated ground truth
(Figure~\ref{fig:ablation}). Three regimes emerged: at
already-resolvable loci (baseline PIP $\geq 0.01$), more coverage
has no visible effect (the causal stays at rank 2);
\textbf{at borderline loci (baseline PIP $\sim 10^{-3}$), more
coverage is dramatically helpful} --- BG1's causal-variant rank
improves from 154 at 10\% coverage to 4 at 100\% (a $38.6\times$
gain) and TGW6's improves from 65 to 3 ($21.6\times$). At
underpowered loci (baseline PIP $\leq 10^{-4}$), more coverage
\emph{degrades} the causal's rank, because added annotation mass is
dispersed across non-causal variants. \emph{Heterogeneity}, not
density, drove the response: rice's hand-curated 269-gene pathway
labels and RicePPINet edges produced cross-seed variance in HBP
output (rank standard deviation 7--18 at borderline loci),
whereas the uniform human cache produced exactly zero variance
across seeds.

\paragraph{Direct intervention.} If the heterogeneity-not-density
principle is causal, then injecting per-variant heterogeneous
prior information into the cache --- without changing GWAS
sumstats, LD, or the fine-mapping algorithm --- should turn the
human cache from inert (HBP $=$ L1) into informative (HBP narrows
beyond L1). We extended the cache schema with an optional
per-variant continuous score that HBP and L1 mix into the
final posterior (Methods). For human, the smallest available
heterogeneous layer was an ENCODE cCRE class flag\cite{encode2020expanded}
(926{,}535 elements across nine classes:
proximal/distal enhancer-like, promoter-like, CTCF-bound, etc.;
$\approx$10--15\% of GTEx-annotated variants gain a class-weighted
score). For Arabidopsis we used TAIR10 CDS-vs-gene-body class;
for rice, snpEff impact class (HIGH/MODERATE/LOW/MODIFIER); for
yeast, BIOGRID physical interactions plus ORF feature class.

\textbf{The intervention transforms the human result and confirms
the principle in three of four species} (Table~\ref{tab:intervention}).
On the same 321 human Pan-UKB leads, adding the cCRE class flag
alone takes HBP narrower than L1 on \textbf{282 of 321 leads (88\%)},
up from 0 of 321 (0\%) without it --- a 88-percentage-point
intervention effect on a fixed locus set. The same direction holds
for Arabidopsis (54\%$\to$94\%) and rice (19\%$\to$78\%); yeast was
already saturated at 100\%. Across all 692 leads, HBP narrows
beyond L1 on 287 (41\%) before the intervention and 634 (92\%)
after.

\begin{figure}[htbp]
\centering
\includegraphics[width=0.85\textwidth]{fig_ablation_rice.pdf}
\caption{\textbf{Multi-omics coverage ablation on real rice GWAS
data.} For each of five curated grain-quality causal loci, HBP
was re-run with 10--100\% of multi-omics edges randomly retained
(10 bootstrap seeds per non-100\% fraction). The rank of the true
causal variant (log scale, lower is better) is plotted against
retained coverage. Three regimes emerge: \emph{resolvable} loci
(green, dashed) are already confident from sumstats and LD alone
and are coverage-insensitive; \emph{borderline} loci (blue, solid)
improve 20--40$\times$ in rank as coverage grows --- the regime
where the relational prior adds the most value;
\emph{underpowered} loci (red, dotted) cannot be rescued by
annotation and degrade slightly under dense priors.}
\label{fig:ablation}
\end{figure}

\begin{table}[htbp]
\centering
\small
\caption{\textbf{Heterogeneity-not-density verified by direct
intervention.} Adding a per-variant continuous prior to the
multi-omics cache (without altering GWAS sumstats, LD references, or
the fine-mapping algorithm) consistently raises the rate at which HBP
narrows the credible set beyond flat-prior L1. The human row is the
cleanest demonstration: a uniform GTEx + STRING cache leaves HBP
exactly equal to L1 on every one of 321 leads, but adding a
class-weighted ENCODE cCRE flag flips that rate from 0\% to 88\% on
the same lead set.}
\label{tab:intervention}
\begin{tabular}{lrrrr}
\toprule
\textbf{Species} & \textbf{Loci} & \textbf{Baseline cache} & \textbf{+heterogeneous prior} & \textbf{$\Delta$} \\
                 &               & \textbf{HBP-tighter}     & \textbf{HBP-tighter}             &                   \\
\midrule
Human       & 321 & 0 / 321 (0\%)   & 282 / 321 (88\%) & \textbf{$+$88\,pp} \\
Arabidopsis &  54 & 29 / 54 (54\%)  & 51 / 54 (94\%)   & $+$41\,pp \\
Rice        &  72 & 14 / 72 (19\%)  & 56 / 72 (78\%)   & $+$58\,pp \\
Yeast       & 245 & 244 / 245 (100\%) & 245 / 245 (100\%) & saturated \\
\midrule
\textbf{Total} & 692 & \textbf{287 (41\%)} & \textbf{634 (92\%)} & \textbf{$+$50\,pp} \\
\bottomrule
\end{tabular}
\end{table}

These intervention results upgrade the heterogeneity-not-density
principle from a cross-species observation to a causal claim:
density alone is insufficient even at $10^7$-variant scale; the
operative quantity is per-variant discriminative information, and
even a coarse nine-class regulatory-element flag is enough to
realise it. Continuous-score replacements --- combined-annotation
PHRED scores, allele-specific transcription-factor-binding
$\Delta$-PWM, deep-learning chromatin-effect predictions, and
species-specific cistromes --- are a natural next direction.

\subsection*{Method selection and extensions}

Different questions call for different methods (Figure~\ref{fig:selection};
full decision table in Supplementary Table~S3). All GraphGWAS methods share a unified
input interface---graph database, BGEN file, or Pan-UKB-style
summary statistics---and a unified output:
\texttt{AssociationResult} and \texttt{CredibleSet} typed graph
nodes queryable by the database's native query language alongside
the biological annotations that produced them. This is the
practical realisation of the graph-native claim: the fine-mapping
output is itself a graph object, not a decorated table, and the
downstream question ``which credible-set variants are in a
druggable pathway?'' is a single traversal rather than a multi-step
enrichment cascade.

The same graph substrate extends naturally beyond single-variant
fine-mapping. As a brief preview of forthcoming work, LD-pruned
co-occurrence (M1) reformulates pairwise epistasis discovery as a
search over a maximal $r^2 < \tau$ independent set of the LD graph,
testing only pairs within it. Theorem~1 (Supplementary Note~S1)
proves the search space shrinks to $\approx 1/k(\tau)^2$ of the
exhaustive $M(M-1)/2$ pairs; on 1000 Genomes chromosome~22
($M = 102{,}467$, $\tau = 0.5$), $k(\tau) \approx 0.003$ yields a
\textbf{42{,}000$\times$ search-space reduction}
($5.2 \times 10^9 \rightarrow \approx 31$~K candidate pairs) with
no loss of ground-truth interaction detection
(Supplementary Fig.~S2). Full epistasis benchmarking is the subject
of a companion manuscript in preparation.

\begin{figure}[htbp]
\centering
\includegraphics[width=0.85\textwidth]{fig7_method_selection.pdf}
\caption{\textbf{Method selection decision tree.} Recommended
GraphGWAS method by signal strength, annotation informativeness,
LD complexity, and whether the workload is per-locus or
genome-wide. For strong signal without annotations,
SuSiE~/~FINEMAP remain the standard. For strong signal with
eQTL annotations, HBP matches accuracy at 20--30$\times$ the
speed. For weak signal with informative annotations, L1 wins
decisively (27--2 over SuSiE; Figure~\ref{fig:weak_signal}). For
dense LD with annotations, L1's graph prior breaks LD ties that
flat-prior methods cannot. For genome-wide workloads, SBayesRC is
appropriate; HBP/L1 target high-precision per-locus refinement at
SBayesRC-surfaced leads. All methods are accessible through the
same input interface (graph database, BGEN, or Pan-UKB-style
sumstats) and return \texttt{CredibleSet} nodes queryable
alongside biological annotations.}
\label{fig:selection}
\end{figure}

